% Generated draft: requires mathematical and source review.
\documentclass[11pt]{article}
\usepackage[margin=25mm]{geometry}
\usepackage{fontspec}
\setmainfont{Latin Modern Roman}
\usepackage{amsmath,amssymb,mathtools,unicode-math}
\setmathfont{Latin Modern Math}
\usepackage{enumitem,needspace,xurl,hyperref}
\setlength{\parindent}{0pt}
\setlength{\parskip}{5pt}
\setlength{\emergencystretch}{3em}
\newcommand{\N}{\mathbb{N}}
\newcommand{\Z}{\mathbb{Z}}
\newcommand{\Q}{\mathbb{Q}}
\newcommand{\R}{\mathbb{R}}
\newcommand{\C}{\mathbb{C}}
\newcommand{\F}{\mathbb{F}}
\newcommand{\A}{\mathbb{A}}
\newcommand{\PP}{\mathbb{P}}
\newcommand{\E}{\mathbb{E}}
\newcommand{\eps}{\varepsilon}
\newcommand{\dd}{\,\mathrm{d}}
\begin{document}
\section*{UnsolvedMath ID: 20000817}
\textbf{Source problem code:} catalogue entry
\subsection{Definitions and mathematical statement}
Fix coordinates on $\A_k^n$ and an irreducible component $C$ of $\operatorname{Hilb}^d(\A_k^n)$. Its full monomial signature is
\[
 \mathcal M(C)=\{I\subset k[x_1,\ldots,x_n]:I\text{ monomial},\ \dim_k k[x_1,\ldots,x_n]/I=d,\ [I]\in C\}.
\]
For $\operatorname{Hilb}^P(\PP_k^n)$, use saturated homogeneous monomial ideals defining the corresponding projective subschemes. The coordinates are fixed; component means an irreducible support, not an additional nonreduced scheme structure to reconstruct.

\textbf{Problem 16.} Does the set of monomial ideals in a component determine that component? Equivalently, must $\mathcal M(C)=\mathcal M(C')$ imply $C=C'$? Ask first for affine schemes of points and more generally for an arbitrary admissible projective Hilbert polynomial $P$. Equality of complete signatures is the issue, rather than whether distinct components can share a monomial point.

\textbf{Reported characteristic-zero results.} Over an algebraically closed field of characteristic zero, the aid states that the signature map is injective on all components of $\operatorname{Hilb}^d(\A^n)$ for every $1\leq d\leq8$ and every $n\geq1$.

More generally, any component containing a monomial point smooth in the ambient Hilbert scheme is determined by its full signature: such a smooth point belongs to a unique irreducible component. The aid applies this to the smoothable component, the closure of the locus of $d$ distinct reduced points, using the curvilinear ideal
\[
 I_{\mathrm{curv}}=(x_1^d,x_2,\ldots,x_n).
\]
Thus that component is always distinguished in the reported setting.

The aid also describes a singular monomial ideal of length eight lying on both components of $\operatorname{Hilb}^8(\A^4)$. Its explicit generators are absent from the supplied summary. This shared point does not contradict injectivity of full signatures: the two complete sets remain different.

For projective Hilbert schemes the summary reports the analogous distinction for the lexicographic component and, in characteristic zero, for all components whenever there are at most two Borel-fixed ideals. A Borel-fixed ideal is fixed by the standard upper-triangular coordinate group; the exact theorem hypotheses and references are not supplied. No conclusion for all projective polynomials, arbitrary affine lengths or positive characteristic follows from these cases.

The original nested remark tentatively says Richard Liebling might have given a counterexample for the analogous curve-Hilbert-scheme question in his thesis. This uncertainty is retained: the record supplies no thesis identification or verified example, so it is not declared a settled counterexample.
\subsection{Short English statement}
Can the complete set of monomial points identify a Hilbert-scheme component? The aid reports yes for affine lengths through eight and gives a smooth-monomial-point criterion, including the smoothable component. Components may still share individual monomial points. The broader projective question and tentative thesis counterexample remain qualified.
\subsection{Sources}
\begin{itemize}
\item Local UnsolvedMath registry, permanent ID $20000817$, \texttt{catalogue entry}; canonical \texttt{aim-arithmetic-geometry-notes.json}, \texttt{notes[62]}.
\item AIM workshop ``Components of Hilbert schemes,'' PDF Problem 16: \url{https://aimath.org/WWN/hilbertschemes/hilbertschemes.pdf}. Original remark tentatively mentions Richard Liebling's thesis; its trailing page-header text is an extraction artifact.
\item Research attempt 2, \texttt{partial\_result}, run status \texttt{open}, internal rigor \texttt{proved}, low novelty confidence; current registry partially solved. Explicitly machine-generated, not peer reviewed, and to be independently verified before citation. Original literature field empty; no complete references for the low-length or projective results are supplied.
\end{itemize}
\end{document}
